\section{But de l'annexe}\label{annexe-energie-but}

Cette annexe détaille la preuve utilisée pour contrôler les variables
\(H\) et \(V\) dans le noyau hybride. La difficulté vient du fait qu'une
même fonctionnelle ne possède pas les deux propriétés dont nous avons
besoin :

\begin{itemize}
\item
  la fonctionnelle \(\mathcal L\), qui contient un terme croisé, permet
  d'obtenir une décroissance pendant le flot ;
\item
  l'énergie \(E\), sans terme croisé, ne peut pas augmenter lors d'un
  saut.
\end{itemize}

La preuve combine ces deux fonctionnelles. Elle donne une condition
suffisante de contrôle global des trajectoires. Elle ne couvre pas le
jeu de paramètres principal, notamment \(\lambda=2\).

\section{Notations et hypothèses}\label{annexe-energie-hypotheses}

Pour alléger les écritures, on rassemble les variables nodales dans les
vecteurs \(H=(H_1,\ldots,H_N)\) et \(V=(V_1,\ldots,V_N)\), munis de la
norme euclidienne produit :

\[
\|H\|^2=\sum_{i=1}^N\|H_i\|^2,
\qquad
\|V\|^2=\sum_{i=1}^N\|V_i\|^2.
\]

Entre deux événements,

\[
\dot H_i=V_i,
\]

\[
\dot V_i
=
-\Gamma(\tau_i)V_i
-M(\tau_i)H_i
+f_i(t),
\]

où

\[
f_i(t)
=
\zeta\,\tanh\!\bigl((\widetilde A H)_i\bigr)+u_i(t).
\]

Comme \(\tanh\) est bornée et \(\|u_i(t)\|\leq u_{\max}\),

\[
\|f_i(t)\|
\leq
B,
\qquad
B:=\zeta\sqrt d+u_{\max}.
\]

Les bornes déjà établies sur la trace donnent

\[
m_0I_d\preceq M(\tau_i)\preceq I_d,
\qquad
m_0=\frac{1}{1+\lambda},
\]

\[
\gamma_0
\leq
\Gamma(\tau_i)
\leq
\bar\Gamma,
\]

et

\[
\|\dot M(\tau_i)\|
\leq
\mu.
\]

Ici, avec la même borne de trace que dans l'Annexe A,
\[
\bar\tau
=
\max_i\max\left\{\|\tau_i(t_0)\|_2,
\frac{\alpha\sqrt d}{\beta}\right\},
\]
et

\[
\bar\Gamma
=
\gamma_0+\gamma_1\bar\tau,
\qquad
\mu
=
L_M(\lambda)
\bigl(\alpha\sqrt d+\beta\bar\tau\bigr).
\]

\section{Les deux fonctionnelles}\label{annexe-energie-deux-fonctionnelles}

On définit d'abord l'énergie sans terme croisé :

\[
E
=
\sum_{i=1}^N
\left[
\frac12\|V_i\|^2
+
\frac12H_i^\top M(\tau_i)H_i
\right].
\]

On introduit ensuite

\[
\mathcal L
=
E
+
\varepsilon
\sum_{i=1}^N H_i\cdot V_i,
\qquad
\varepsilon>0.
\]

Le terme croisé permet de faire apparaître, dans la dérivée de
\(\mathcal L\), un terme négatif en \(\|H\|^2\). Il complique en
revanche l'analyse des sauts.

\section{Équivalence entre \texorpdfstring{$E$}{E} et \texorpdfstring{$\mathcal L$}{L}}
\label{annexe-energie-equivalence}

La borne spectrale \(M(\tau_i)\succeq m_0I_d\) donne

\[
E
\geq
\frac12\|V\|^2
+
\frac{m_0}{2}\|H\|^2.
\]

Par conséquent,

\[
\|V\|
\leq
\sqrt{2E},
\qquad
\|H\|
\leq
\sqrt{\frac{2E}{m_0}}.
\]

L'inégalité de Cauchy--Schwarz donne alors

\[
\left|
\sum_{i=1}^N H_i\cdot V_i
\right|
\leq
\|H\|\,\|V\|
\leq
\frac{2E}{\sqrt{m_0}}.
\]

Ainsi,

\[
\mathcal L
\geq
E-\frac{2\varepsilon}{\sqrt{m_0}}E
=
c^{-}\,E,
\]

et

\[
\mathcal L
\leq
E+\frac{2\varepsilon}{\sqrt{m_0}}E
=
c^{+}\,E,
\]

avec

\[
c^{-}
=
1-\frac{2\varepsilon}{\sqrt{m_0}},
\qquad
c^{+}
=
1+\frac{2\varepsilon}{\sqrt{m_0}}.
\]

Si

\[
0<\varepsilon<\frac{\sqrt{m_0}}{2},
\]

alors \(c^{-}>0\) et

\[
\boxed{
c^{-}\,E
\leq
\mathcal L
\leq
c^{+}\,E
}.
\]

Le symbole \(c^{-}\,E\) désigne le produit de la constante \(c^{-}\) par
l'énergie \(E\).

\section{Dérivée de \texorpdfstring{$\mathcal L$}{L} pendant le flot}
\label{annexe-energie-derivee-flot}

Nous calculons la dérivée pour un nœud, puis nous sommons.

D'une part,

\[
\frac{\mathrm d}{\mathrm dt}
\left(\frac12\|V_i\|^2\right)
=
V_i\cdot\dot V_i
=
-\Gamma_i\|V_i\|^2
-
V_i^\top M_iH_i
+
V_i\cdot f_i,
\]

où, pour alléger l'écriture,

\[
\Gamma_i=\Gamma(\tau_i),
\qquad
M_i=M(\tau_i).
\]

D'autre part,

\[
\frac{\mathrm d}{\mathrm dt}
\left(
\frac12H_i^\top M_iH_i
\right)
=
V_i^\top M_iH_i
+
\frac12H_i^\top\dot M_iH_i.
\]

Les deux termes mixtes \(V_i^\top M_iH_i\) s'annulent.

Enfin,

\[
\frac{\mathrm d}{\mathrm dt}
\left(
\varepsilon H_i\cdot V_i
\right)
=
\varepsilon\|V_i\|^2
+
\varepsilon H_i\cdot\dot V_i,
\]

donc

\[
\frac{\mathrm d}{\mathrm dt}
\left(
\varepsilon H_i\cdot V_i
\right)
=
\varepsilon\|V_i\|^2
-
\varepsilon\Gamma_iH_i\cdot V_i
-
\varepsilon H_i^\top M_iH_i
+
\varepsilon H_i\cdot f_i.
\]

En additionnant,

\[
\begin{aligned}
\dot{\mathcal L}
={}&
-\sum_i(\Gamma_i-\varepsilon)\|V_i\|^2
-\varepsilon\sum_iH_i^\top M_iH_i
-\varepsilon\sum_i\Gamma_iH_i\cdot V_i
\\
&+
\frac12\sum_iH_i^\top\dot M_iH_i
+
\sum_iV_i\cdot f_i
+
\varepsilon\sum_iH_i\cdot f_i.
\end{aligned}
\]

\section{Contrôle de la partie quadratique}
\label{annexe-energie-partie-quadratique}

Les bornes uniformes donnent

\[
-\sum_i(\Gamma_i-\varepsilon)\|V_i\|^2
\leq
-(\gamma_0-\varepsilon)\|V\|^2,
\]

\[
-\varepsilon\sum_iH_i^\top M_iH_i
\leq
-\varepsilon m_0\|H\|^2,
\]

et

\[
\frac12\sum_iH_i^\top\dot M_iH_i
\leq
\frac{\mu}{2}\|H\|^2.
\]

Pour le terme croisé, on utilise

\[
\varepsilon\bar\Gamma\|H\|\,\|V\|
\leq
\frac{\varepsilon m_0}{4}\|H\|^2
+
\frac{\varepsilon\bar\Gamma^2}{m_0}\|V\|^2.
\]

On obtient donc

\[
\dot{\mathcal L}
\leq
-A_V\|V\|^2
-A_H\|H\|^2
+
\sum_iV_i\cdot f_i
+
\varepsilon\sum_iH_i\cdot f_i,
\]

avec

\[
A_V
=
\gamma_0
-
\varepsilon
\left(
1+\frac{\bar\Gamma^2}{m_0}
\right),
\]

et

\[
A_H
=
\frac{3\varepsilon m_0}{4}
-
\frac{\mu}{2}.
\]

Pour avoir \(A_V>0\) et \(A_H>0\), il suffit de choisir
\(\varepsilon\) tel que

\[
\frac{2\mu}{3m_0}
<
\varepsilon
<
\frac{\gamma_0}
{2\left(1+\bar\Gamma^2/m_0\right)}.
\]

Un tel choix est possible si

\[
\boxed{
\mu
<
\frac{3m_0\gamma_0}
{4\left(1+\bar\Gamma^2/m_0\right)}
}
\tag{S}
\]

Cette condition est suffisante, non nécessaire. Les constantes utilisées ici sont uniformes sur toute la plage admissible de \(\|\tau\|\) et ne tiennent pas compte de leurs dépendances communes en \(\tau\). Leur conservatisme quantitatif est discuté à la section~\ref{annexe-energie-limites}. Le facteur \(1/2\) dans la borne supérieure imposée à \(\varepsilon\) garde une marge explicite :

\[
A_V\geq\frac{\gamma_0}{2}.
\]

\section{Absorption de l'entrée bornée}
\label{annexe-energie-entree-bornee}

Comme \(\|f_i\|\leq B\),

\[
\sum_iV_i\cdot f_i
\leq
B\sum_i\|V_i\|,
\]

et

\[
\varepsilon\sum_iH_i\cdot f_i
\leq
\varepsilon B\sum_i\|H_i\|.
\]

Pour chaque nœud, l'inégalité de Young donne

\[
B\|V_i\|
\leq
\frac{A_V}{2}\|V_i\|^2
+
\frac{B^2}{2A_V},
\]

et

\[
\varepsilon B\|H_i\|
\leq
\frac{A_H}{2}\|H_i\|^2
+
\frac{\varepsilon^2B^2}{2A_H}.
\]

Après sommation,

\[
\dot{\mathcal L}
\leq
-\frac{A_V}{2}\|V\|^2
-\frac{A_H}{2}\|H\|^2
+b,
\]

où

\[
b
=
\frac{NB^2}{2}
\left(
\frac{1}{A_V}
+
\frac{\varepsilon^2}{A_H}
\right).
\]

Posons

\[
d_0=\min(A_V,A_H).
\]

Comme \(M_i\preceq I_d\),

\[
E
\leq
\frac12
\left(
\|V\|^2+\|H\|^2
\right),
\]

puis

\[
\mathcal L
\leq
\frac{c^{+}}{2}
\left(
\|V\|^2+\|H\|^2
\right).
\]

Il en résulte

\[
\|V\|^2+\|H\|^2
\geq
\frac{2}{c^{+}}\mathcal L.
\]

Par conséquent,

\[
\dot{\mathcal L}
\leq
-\frac{d_0}{c^{+}}\mathcal L+b.
\]

En posant

\[
a=\frac{d_0}{2c^{+}},
\]

on obtient la forme utilisée dans le corps du texte :

\[
\boxed{
\dot{\mathcal L}
\leq
-2a\mathcal L+b
}.
\]

Sur un intervalle de flot de durée \(T\),

\[
\boxed{
\mathcal L(t+T)
\leq
e^{-2aT}\mathcal L(t)
+
\frac{b}{2a}
\left(
1-e^{-2aT}
\right)
}.
\]

En particulier,

\[
\mathcal L(t+T)
\leq
e^{-2aT}\mathcal L(t)
+
\frac{b}{2a}.
\]

\section{Effet exact d'un saut sur les deux fonctionnelles}
\label{annexe-energie-effet-saut}

Lors d'un saut du nœud \(i^\ast\),

\[
V_{i^\ast}^+=r_VV_{i^\ast}^-,
\]

tandis que \(H\), \(\tau\) et \(M(\tau)\) ne changent pas.

Pour l'énergie \(E\),

\[
E^+
=
E^-
-
\frac12(1-r_V^2)
\|V_{i^\ast}^-\|^2.
\]

Ainsi,

\[
\boxed{E^+\leq E^-}.
\]

Pour \(\mathcal L\),

\[
\Delta\mathcal L
=
\mathcal L^+-\mathcal L^-
=
-\frac12(1-r_V)(1+r_V)\|V\|^2
-
\varepsilon(1-r_V)H\cdot V,
\]

où \(H=H_{i^\ast}^-\) et \(V=V_{i^\ast}^-\).

Le deuxième terme peut être positif lorsque \(H\cdot V<0\). On majore
donc

\[
-H\cdot V
\leq
\|H\|\,\|V\|.
\]

Il vient

\[
\Delta\mathcal L
\leq
(1-r_V)
\left[
\varepsilon\|H\|\,\|V\|
-
\frac{1+r_V}{2}\|V\|^2
\right].
\]

Pour \(H\) fixé, le membre de droite est un polynôme concave en
\(y=\|V\|\). Son maximum est atteint pour

\[
y
=
\frac{\varepsilon\|H\|}{1+r_V}.
\]

On obtient

\[
\Delta\mathcal L
\leq
\frac{1-r_V}{2(1+r_V)}
\varepsilon^2\|H\|^2
\leq
\frac{\varepsilon^2}{2}\|H\|^2.
\]

L'équivalence \(c^{-}E\leq\mathcal L\) donne

\[
\|H_{i^\ast}\|^2
\leq
\frac{2E}{m_0}
\leq
\frac{2\mathcal L^-}{c^{-}m_0}.
\]

Par conséquent,

\[
\mathcal L^+
\leq
\left(
1+
\frac{\varepsilon^2}{c^{-}m_0}
\right)
\mathcal L^-.
\]

En posant

\[
q
=
\frac{\varepsilon^2}{c^{-}m_0},
\]

on obtient

\[
\boxed{
\mathcal L^+
\leq
(1+q)\mathcal L^-
}.
\]

Cette borne est volontairement uniforme en \(r_V\). La borne plus
précise conserve le facteur \((1-r_V)/(1+r_V)\).

\section{Temps minimal entre deux sauts d'un même nœud}
\label{annexe-energie-temps-sejour}

Supposons que, sur une portion d'exécution,

\[
\mathcal L\leq K.
\]

Comme \(c^{-}E\leq\mathcal L\),

\[
\|H_i\|
\leq
\sqrt{\frac{2K}{c^{-}m_0}},
\]

et

\[
\|V_i\|
\leq
\sqrt{\frac{2K}{c^{-}}}.
\]

La pression vérifie

\[
\dot P_i
=
\kappa_1\|H_i\|
+
\kappa_2\|V_i\|
-
\kappa_3P_i.
\]

En supprimant le terme négatif,

\[
\dot P_i
\leq
\kappa_1
\sqrt{\frac{2K}{c^{-}m_0}}
+
\kappa_2
\sqrt{\frac{2K}{c^{-}}}
=:\nu(K).
\]

Après un saut,

\[
P_i^+=r_P\Theta.
\]

Avant de pouvoir sauter à nouveau, le même nœud doit donc accumuler au
moins

\[
(1-r_P)\Theta.
\]

Le temps continu nécessaire est au moins

\[
\boxed{
\delta(K)
=
\frac{(1-r_P)\Theta}{\nu(K)}
}.
\]

Dans toute fenêtre semi-ouverte de temps continu
\([t,t+\delta(K))\), chaque nœud saute au plus une fois. Il y a donc au
plus \(N\) sauts dans une telle fenêtre.

\section{Combinaison du flot et des sauts}
\label{annexe-energie-combinaison}

Considérons un intervalle de temps continu de durée \(T\) contenant
\(k\) sauts. La combinaison de l'estimation de flot et de la borne de
saut donne

\[
\mathcal L_{\text{fin}}
\leq
(1+q)^k
\left[
e^{-2aT}\mathcal L_{\text{début}}
+
\frac{b}{2a}
\right].
\]

Cette borne est obtenue en appliquant successivement l'estimation
affine pendant chaque phase de flot et la multiplication par \(1+q\)
à chaque saut. Le terme constant \(b/(2a)\) est lui aussi majoré par le
facteur total \((1+q)^k\).

Une difficulté doit être traitée explicitement : la contraction sur
une fenêtre complète n'exclut pas, à elle seule, une augmentation
transitoire à l'intérieur de cette fenêtre. Nous introduisons donc
deux niveaux :

\begin{itemize}
\item
  \(R\), borne visée aux extrémités des fenêtres ;
\item
  \(\widehat R\), enveloppe autorisée à l'intérieur des fenêtres.
\end{itemize}

Posons

\[
r_J=(1+q)^N,
\qquad
C=\frac{b}{2a},
\]

et

\[
\boxed{
\widehat R
=
r_J(R+C)
}.
\]

Supposons

\[
\mathcal L(0)\leq R.
\]

Tant que \(\mathcal L\leq\widehat R\), une fenêtre de longueur
\(\delta(\widehat R)\) contient au plus \(N\) sauts. Sur toute portion
d'une telle fenêtre,

\[
\mathcal L(t)
\leq
r_J\bigl(R+C\bigr)
=
\widehat R.
\]

Un argument de premier franchissement montre donc que
\(\mathcal L\) ne peut pas dépasser \(\widehat R\) à l'intérieur de la
fenêtre.

À l'extrémité d'une fenêtre complète,

\[
\mathcal L_{\text{suivante}}
\leq
r_J
\left[
e^{-2a\delta(\widehat R)}R+C
\right].
\]

La quantité \(\widehat R=r_J(R+C)\) est d'abord déterminée par le
choix de \(R\), puis \(\delta(\widehat R)\) s'en déduit. La condition
suivante est donc une contrainte sur les paramètres vérifiée a
posteriori, et non une équation implicite en \(R\).

Il suffit alors d'imposer

\[
\boxed{
r_J
e^{-2a\delta(\widehat R)}
\leq
\frac12
}
\]

et

\[
\boxed{
R
\geq
2r_JC
=
r_J\frac{b}{a}
}.
\tag{$Z_R^\sharp$}
\]

Dans ce cas,

\[
\mathcal L_{\text{suivante}}
\leq
\frac{R}{2}
+
r_JC
\leq R.
\]

L'argument peut être répété sur toutes les fenêtres successives.

\section{Résultat final}
\label{annexe-energie-resultat-final}

\begin{quote}
\resultheading{Proposition}{prop:B-bornitude-conditionnelle}{}
Supposons la condition \((S)\). Choisissons
\(\varepsilon<\sqrt{m_0}/2\) dans l'intervalle indiqué plus haut, et
définissons \(a\), \(b\), \(c^{-}\), \(q\), \(r_J\), \(R\) et
\(\widehat R\) comme précédemment. Si
\(\mathcal L(0)\leq R\) et si \((Z_R^\sharp)\) est satisfaite, alors
\[
\mathcal L(t)\leq\widehat R
\]
pour tout temps continu \(t\). Les variables \(H_i\) et \(V_i\)
restent donc bornées le long de toute l'exécution hybride.
\end{quote}

En effet,

\[
\|H_i(t)\|
\leq
\sqrt{\frac{2\widehat R}{c^{-}m_0}},
\qquad
\|V_i(t)\|
\leq
\sqrt{\frac{2\widehat R}{c^{-}}}.
\]

La preuve donne une borne explicite, mais conservative.

\section{Différence avec la formulation abrégée du corps du texte}
\label{annexe-energie-correction-formulation}

Une formulation antérieure utilisait directement \(R\) à la fois comme
borne aux extrémités des fenêtres et comme borne à tout instant. Cette
présentation était trop rapide : un saut peut augmenter
\(\mathcal L\) avant que le flot ne compense cette augmentation.

L'enveloppe \(\widehat R\) sépare désormais :

\[
\text{contrôle aux extrémités des fenêtres}
\quad\text{et}\quad
\text{augmentation transitoire possible}.
\]

Le résultat nécessaire au manuscrit reste le même : les trajectoires
restent dans un ensemble borné sous des conditions suffisantes
explicites. La borne globale est \(\widehat R\), et non \(R\).

\section{Portée et limites}\label{annexe-energie-limites}

La preuve repose sur deux conditions suffisantes :

\[
(S)
\qquad\text{et}\qquad
(Z_R^\sharp).
\]

Elles sont conservatrices pour plusieurs raisons :

\begin{itemize}
\item
la constante \(L_M(\lambda)\) utilisée pour borner la différentielle de \(M\) est elle-même une majoration globale conservative ;
\item
les bornes sur \(\dot M\), la coercivité de \(M\) et la friction \(\Gamma\) sont uniformes et traitées indépendamment, alors qu'elles dépendent toutes de la même variable \(s=\|\tau\|\) ;
\item
le terme croisé est absorbé par une inégalité de Young avec un choix uniforme sur toute la plage admissible de \(s\) ;
\item
le contrôle du saut oublie le facteur favorable
\((1-r_V)/(1+r_V)\) ;
\item
tous les sauts possibles d'une fenêtre sont majorés par le même
facteur \(1+q\) ;
\item
aucune corrélation entre la phase du flot et les instants de saut
n'est utilisée.
\end{itemize}

Pour les paramètres de référence de la réduction bidimensionnelle étudiée numériquement, en particulier \(\lambda=2\), la condition \((S)\) n'est pas satisfaite. Son écart important ne doit toutefois pas être interprété comme une indication d'instabilité du régime considéré.

À \(d=2\), la condition \((S)\) telle qu'elle résulte des constantes uniformes utilisées ci-dessus est dépassée d'un facteur d'environ \(5{,}5\times10^3\). Le remplacement de la borne globale \(L_M(\lambda)\) par le supremum exact de \(\|DM\|\) ramène ce facteur à environ \(2{,}2\times10^3\). Une réévaluation ponctuelle des mêmes termes en fonction de \(s=\|\tau\|\), utilisant notamment la coercivité locale de \(M\) et la friction locale, ramène encore ce déficit à environ \(1{,}027\). Cette optimisation choisit toutefois \(\varepsilon\) indépendamment pour chaque valeur de \(s\) : elle fournit une minoration du déficit accessible à cette famille de majorations, mais ne définit pas une fonction de Lyapunov unique admissible.

Ces calculs ne constituent pas une nouvelle preuve de bornitude. Ils montrent seulement que l'essentiel de l'écart de \((S)\) provient des majorations uniformes utilisées dans sa dérivation. À titre de contrôle complémentaire, une fonction quadratique commune beaucoup plus simple, avec potentiel isotrope constant et terme croisé constant, ramène le déficit à environ \(1{,}12\) pour \(d=2\), sans toutefois fermer la condition.

Les simulations indiquent que les trajectoires restent bornées, mais cette annexe ne le démontre pas pour le régime principal.

Le résultat établi reste donc :

$$
\boxed{
\text{contrôle global démontré sous }(S)\land(Z_R^\sharp),
\quad
\text{question ouverte dans le régime principal}.
}
$$

